% =====================================================================
% ARTIGO FINAL — Detecção e Mitigação Inteligente de Anomalias em SDN
% com LLMs Open Source: correlação Projeto de Pesquisa (goal.pdf) x Sistema
% Implementado (NOC Lite).
% Autor: George (escrita acadêmica ABNT). Codificação: UTF-8 / pdflatex.
% Compilação: pdflatex artigo.tex 2x (TinyTeX; abntex2cite para citação
% autor-data ABNT NBR 10520). Necessita do pacote setspace (instalado via tlmgr).
% =====================================================================
\documentclass[12pt,a4paper]{article}

% --- Codificação e idioma (PT-BR). TinyTeX minimal: acentos explícitos via
% \~a, \c{c}, \'a; nao se usa opcao de linguagem do babel para garantir
% compilacao em qualquer instalacao. ---
\usepackage[utf8]{inputenc}
\usepackage[T1]{fontenc}
\usepackage{indentfirst}      % indentar primeiro parágrafo (norma ABNT)

% --- Citacao ABNT autor-data (NBR 10520) via abntex2cite ---

% --- Utilitários ---
\usepackage{graphicx}
\usepackage{hyperref}
\usepackage{float}
\usepackage{longtable}
\usepackage{booktabs}
\usepackage{amsmath}
\usepackage{caption}

% --- Hiperlinks em preto (evita caixas coloridas na impressão) ---
\hypersetup{colorlinks=true, linkcolor=black, citecolor=black,
            urlcolor=black, pdfborder={0 0 0}}

\graphicspath{{../figures/}{figuras/}}

% ---------------------------------------------------------------------
% Comandos de citação (autor, ano) — gerados a partir das 6 áreas.
% Cada \rXX aponta para um \bibitem em referencias.tex. Não há
% referências inventadas: todas provêm dos arquivos areaNN_*.txt.
% (o \input foi movido para APÓS \begin{document} para evitar
%  "Missing begin document" no preâmbulo)

% ---------------------------------------------------------------------
\title{\textbf{Detecção e Mitigação Inteligente de Anomalias em Redes
Definidas por Software com Grandes Modelos de Linguagem Open Source:
correlação entre o Projeto de Pesquisa e o Sistema Implementado (NOC Lite)}}

\author{Estudo de correlação elaborado a partir do Projeto de Pesquisa
de \mbox{Rodrigues} (2026, Edital n\textsuperscript{o} 26/2026) e do
sistema \textit{NOC Lite} implementado \\[2mm]
\textit{Artigo técnico-científico — revisão integrativa e análise arquitetural}}

\date{Julho de 2026}

% ===== macros autor-data
% =====================================================================
% Macros de citação ABNT autor-data (NBR 10520)
% =====================================================================
\makeatletter
\newcommand{\citea}[1]{\begingroup\edef\@currentlabel{\csname citea@#1\endcsname}\@currentlabel\endgroup}
\newcommand{\citea@rA11}{McKEOWN et al, 2008}
\newcommand{\citea@rA110}{XIA et al, 2015}
\newcommand{\citea@rA114}{MEDVED et al, 2014}
\newcommand{\citea@rA115}{BERDE et al, 2014}
\newcommand{\citea@rA116}{NTT Laboratories. Ryu SDN Framework - Open-source SDN Platform Software. NTT Technical Review, 2014}
\newcommand{\citea@rA12}{GUDE et al, 2008}
\newcommand{\citea@rA122}{PORRAS et al, 2012}
\newcommand{\citea@rA123}{LI et al, 2016}
\newcommand{\citea@rA124}{BHUIYAN et al, 2023}
\newcommand{\citea@rA126}{LATIF et al, 2022}
\newcommand{\citea@rA127}{CHIOSI et al, 2012}
\newcommand{\citea@rA128}{ETSI NFV ISG. GS NFV 002 V1.1.1: Network Functions Virtualisation (NFV); Architectural Framework. ETSI et al, 2013}
\newcommand{\citea@rA129}{ETSI NFV ISG. GS NFV-MAN 001 V1.1.1: Network Functions Virtualisation (NFV); Management and Orchestration. ETSI et al, 2014}
\newcommand{\citea@rA130}{MIJUMBI et al, 2015}
\newcommand{\citea@rA131}{MIJUMBI et al, 2016}
\newcommand{\citea@rA132}{LAGHRISSI et al, 2019}
\newcommand{\citea@rA133}{LEIVADEAS et al, 2022}
\newcommand{\citea@rA134}{MORALES et al, 2017}
\newcommand{\citea@rA135}{VELASCO et al, 2022}
\newcommand{\citea@rA14}{LANTZ et al, 2010}
\newcommand{\citea@rA15}{PFAFF et al, 2015}
\newcommand{\citea@rA19}{KREUTZ et al, 2015}
\newcommand{\citea@rA21}{BRUTLAG, 2000}
\newcommand{\citea@rA210}{BHUYAN et al, 2014}
\newcommand{\citea@rA212}{BUCZAK et al, 2016}
\newcommand{\citea@rA213}{CORTES et al, 1995}
\newcommand{\citea@rA217}{CARUANA et al, 2006}
\newcommand{\citea@rA22}{ROBERTS, 1959}
\newcommand{\citea@rA222}{HODO et al, 2017}
\newcommand{\citea@rA224}{ALDUAILIJ et al, 2022}
\newcommand{\citea@rA225}{NEGANDHI et al, 2016}
\newcommand{\citea@rA226}{WILLEMS et al, 2023}
\newcommand{\citea@rA227}{SHEHAB et al, 2024}
\newcommand{\citea@rA229}{STANIFORD-CHEN et al, 1996}
\newcommand{\citea@rA231}{CANNADY, 1998}
\newcommand{\citea@rA232}{LAKHINA et al, 2004}
\newcommand{\citea@rA25}{MOAYEDI et al, 2008}
\newcommand{\citea@rA27}{FONDO-FERREIRO et al, 2026}
\newcommand{\citea@rA28}{CHANDOLA et al, 2009}
\newcommand{\citea@rA29}{AHMED et al, 2016}
\newcommand{\citea@rA31}{LIU et al, 2025}
\newcommand{\citea@rA310}{TOUVRON et al, 2023}
\newcommand{\citea@rA311}{TOUVRON et al, 2023}
\newcommand{\citea@rA312}{JIANG et al, 2023}
\newcommand{\citea@rA313}{ALMAZROUEI et al, 2023}
\newcommand{\citea@rA314}{BAI et al, 2023}
\newcommand{\citea@rA315}{YANG et al, 2024}
\newcommand{\citea@rA316}{BIGSCIENCE WORKSHOP et al, 2022}
\newcommand{\citea@rA317}{GEMMA TEAM (Google DeepMind). Gemma: Open Models Based on Gemini Research and Technology. arXiv preprint arXiv:2403.08295, 2024}
\newcommand{\citea@rA319}{VASWANI et al, 2017}
\newcommand{\citea@rA32}{LIU et al, 2024}
\newcommand{\citea@rA320}{OUYANG et al, 2022}
\newcommand{\citea@rA321}{WEI et al, 2022}
\newcommand{\citea@rA324}{HOSSAIN et al, 2025}
\newcommand{\citea@rA325}{SDN-Xplain: LLM-Driven Anomaly Detection and Human-Interpretable Explanation in SDN. SKKU IoT Lab / KICS Winter 2026 (preprint), 2026}
\newcommand{\citea@rA327}{Lightweight Fine-Tuning of LLMs for Explainable Intrusion Detection in SDN. arXiv/ResearchGate preprint, 2025}
\newcommand{\citea@rA33}{BOATENG et al, 2024}
\newcommand{\citea@rA332}{ADANZA et al, 2025}
\newcommand{\citea@rA333}{Towards Intent-Based Network Management: Large Language Models for Intent Extraction in 5G Core Networks. IEEE / GenAI Net research library, 2024}
\newcommand{\citea@rA334}{Exploring LLM-Based Agents for Root Cause Analysis. ACM SIGMETRICS / arXiv:2403.04123, 2024}
\newcommand{\citea@rA335}{Empowering AIOps: Leveraging Large Language Models for IT Operations Management. arXiv:2501.12461, 2025}
\newcommand{\citea@rA336}{WANG et al, 2025}
\newcommand{\citea@rA337}{LLM-Driven Multi-Agent Framework for Autonomous Network Slice Management. IEEE, 2026}
\newcommand{\citea@rA338}{CHEN et al, 2024}
\newcommand{\citea@rA34}{HANG et al, 2024}
\newcommand{\citea@rA35}{HUANG et al, 2024}
\newcommand{\citea@rA36}{HONG et al, 2025}
\newcommand{\citea@rA37}{ZHANG et al, 2025}
\newcommand{\citea@rA39}{TANG et al, 2024}
\newcommand{\citea@rA41}{HU et al, 2021}
\newcommand{\citea@rA411}{HU et al, 2023}
\newcommand{\citea@rA412}{ZHANG et al, 2024}
\newcommand{\citea@rA414}{DING et al, 2022}
\newcommand{\citea@rA415}{HAN et al, 2024}
\newcommand{\citea@rA42}{DETTMERS et al, 2023}
\newcommand{\citea@rA422}{BROWN et al, 2005}
\newcommand{\citea@rA424}{ACHIAM et al, 2023}
\newcommand{\citea@rA425}{MENG et al, 2025}
\newcommand{\citea@rA426}{GINIGE et al, 2024}
\newcommand{\citea@rA427}{CUI et al, 2025}
\newcommand{\citea@rA43}{ZHANG et al, 2023}
\newcommand{\citea@rA431}{SUI et al, 2024}
\newcommand{\citea@rA432}{HEGSELMANN et al, 2023}
\newcommand{\citea@rA44}{LIU et al, 2024}
\newcommand{\citea@rA45}{HOULSBY et al, 1902}
\newcommand{\citea@rA49}{LIU et al, 2022}
\newcommand{\citea@rA51}{ESSENCIAL Danilevsky et al, 2010}
\newcommand{\citea@rA511}{ESSENCIAL Ali et al, 2023}
\newcommand{\citea@rA512}{ESSENCIAL Neupane et al, 2022}
\newcommand{\citea@rA517}{RECOMENDADO Sarker, 2024}
\newcommand{\citea@rA52}{ESSENCIAL Arrieta et al, 2019}
\newcommand{\citea@rA520}{COMPLEMENTAR Houssel et al, 2024}
\newcommand{\citea@rA521}{COMPLEMENTAR Pizarro et al, 2025}
\newcommand{\citea@rA522}{COMPLEMENTAR Lightweight Fine-Tuning of LLMs for Explainable Intrusion Detection in SDN. arXiv:2510.23313 (preprint), 2025}
\newcommand{\citea@rA53}{ESSENCIAL Ribeiro et al, 2016}
\newcommand{\citea@rA531}{ESSENCIAL Kadavath et al, 2022}
\newcommand{\citea@rA532}{RECOMENDADO Guo et al, 2017}
\newcommand{\citea@rA533}{RECOMENDADO Lin et al, 2022}
\newcommand{\citea@rA534}{RECOMENDADO Zhu et al, 2023}
\newcommand{\citea@rA537}{ESSENCIAL Wei et al, 2022}
\newcommand{\citea@rA54}{ESSENCIAL Lundberg et al, 2017}
\newcommand{\citea@rA61}{ESSENCIAL - Dataset UNSW-NB15 (descrição oficial / paper fundador) Moustafa et al, 2015}
\newcommand{\citea@rA610}{RECOMENDADO - Dataset InSDN et al, 2020}
\newcommand{\citea@rA614}{RECOMENDADO - Data augmentation em IoT (desbalanceamento) Zhang et al, 2022}
\newcommand{\citea@rA615}{RECOMENDADO - Geração sintética e privacy-preserving de traces Aceto et al, 2024}
\newcommand{\citea@rA616}{ESSENCIAL - AnoLLM et al, 2025}
\newcommand{\citea@rA617}{RECOMENDADO - Survey LLMs em dados tabulares Fang et al, 2024}
\newcommand{\citea@rA618}{RECOMENDADO - Detecção zero-shot de anomalias tabulares com LLMs Li et al, 2024}
\newcommand{\citea@rA62}{ESSENCIAL - Avaliação estatística do UNSW-NB15 Moustafa et al, 2016}
\newcommand{\citea@rA623}{RECOMENDADO - Augmentation GAN para classes minoritárias em rede Kang et al, 2025}
\newcommand{\citea@rA629}{COMPLEMENTAR - Testbed híbrido SDN et al, 2023}
\newcommand{\citea@rA63}{ESSENCIAL - Metodologia de geração CIC (base do CSE-CIC-IDS2018) Sharafaldin et al, 2018}
\newcommand{\citea@rA630}{RECOMENDADO - Survey de geração sintética de tráfego de rede Synthetic Network Traffic Data Generation: A Comparative Analysis. arXiv preprint, 2024}
\newcommand{\citea@rA64}{ESSENCIAL - Dataset CSE-CIC-IDS2018 (descrição oficial) Canadian Institute for Cybersecurity (CIC) \& Communications Security Establishment (CSE). CSE-CIC-IDS, 2018}
\newcommand{\citea@rA65}{ESSENCIAL - Survey de datasets de NIDS Ring et al, 2019}
\newcommand{\citea@rA66}{RECOMENDADO - Análise do KDD Cup 99 / NSL-KDD et al, 2009}
\newcommand{\citea@rA69}{ESSENCIAL - Mininet (emulador para SDN) Lantz et al, 2010}
\newcommand{\citea@rRODRIGUES}{Rodrigues, João Pedro Santos, 2026}
% ===== end macros autor-data


\begin{document}
\maketitle

% =====================================================================
\begin{abstract}
Este artigo correlaciona o projeto de pesquisa que propõe a detecção e a
mitigação inteligente de anomalias em redes definidas por software (SDN) por
meio de grandes modelos de linguagem (LLMs) \textit{open source} ajustados por
PEFT/LoRA, com o sistema efetivamente implementado -- o \textit{NOC Lite}.
Partindo de uma revisão da literatura estruturada em seis eixos (SDN/NFV,
anomalias por ML, LLMs para NetOps, PEFT/LoRA, XAI e datasets), descreve-se a
arquitetura e o pipeline de classificação do sistema construído, composto por
três classificadores (heurístico, árvore de decisão e IA generativa via Google
Gemini). Em seguida, apresenta-se o mapeamento artigo$\times$sistema e a
taxonomia de gaps estruturais -- destacando-se a ausência de SDN, o uso de
modelo fechado em vez de LLM \textit{open source}, a inexistência de
ajuste fino LoRA/QLoRA, a falta de mitigação em ciclo fechado e a dependência
exclusiva de dados sintéticos. Discutem-se as convergências preservadas
(baselines, tokenização, justificativa textual, métricas de latência) e
propõem-se diretrizes de arquitetura para o alinhamento futuro, incluindo um
ciclo fechado com LLM local, LoRA e módulo de mitigação. Conclui-se que o NOC
Lite realiza parcela da intenção do projeto, mas exige evolução arquitetural
substantiva nos eixos centrais para aproximar-se da proposta norteadora.
\end{abstract}

\noindent\textbf{Palavras-chave:} Redes Definidas por Software; Anomalias de Rede;
Grandes Modelos de Linguagem; PEFT/LoRA; Inteligência Artificial Explicável.

\vspace{6pt}
\hrule
\vspace{10pt}

% =====================================================================
\section{Introdução}
\label{sec:intro}

Redes de computadores modernas apresentam alta complexidade, heterogeneidade de
aplicações e demandas rigorosas de desempenho, confiabilidade e segurança
\citea{rRODRIGUES}. Nesse contexto, a identificação e a classificação de fenômenos
anômalos -- congestionamento transiente, falhas de enlace, degradação de
desempenho -- tornam-se atividades fundamentais para a operação eficiente de
sistemas distribuídos \citea{rA39}. Tradicionalmente, a análise de anomalias depende
de intervenção manual e conhecimento especializado, o que a torna propensa a
erros e com baixa agilidade de adaptação em tempo real \citea{rA39}. Modelos baseados
em assinaturas ou regras pré-definidas são frequentemente reativos e possuem
capacidade limitada para detectar ameaças desconhecidas \citea{rA512}.

Nos últimos anos, grandes modelos de linguagem demonstraram capacidades de
interpretação, raciocínio contextual e generalização a partir de descrições
textuais e dados estruturados \citea{rA31}. Diferentemente de classificadores
tradicionais, tais modelos conseguem combinar múltiplas informações, lidar com
incertezas e produzir não apenas uma decisão, mas também uma justificativa
textual \citea{rA321}. O paradigma SDN, ao separar o plano de controle do plano de
dados, fornece a agilidade necessária para que agentes de IA percebam o estado
da rede e executem ações de otimização ou segurança de forma autônoma \citea{rA19}.

O projeto de pesquisa norteador \citea{rRODRIGUES} propõe desenvolver um sistema
autônomo baseado em LLMs \textit{open source} (ex.: Llama, Falcon) e em
ajuste fino eficiente de parâmetros (PEFT), notadamente LoRA/QLoRA, para
detectar anomalias de performance em SDN e mitigá-las via reconfiguração
proativa de topologia virtual (VNT) e roteamento, em ciclo fechado. Contudo, o
sistema implementado -- o \textit{NOC Lite} -- monitora \textit{hosts}
genéricos (endereços IP/DNS) por sondas ICMP/TCP/HTTP/DNS/SNMP, emprega o
modelo fechado Google Gemini como classificador generativo e não realiza
mitigação nem reconfiguração de rede.

\textbf{Problema e gap.} Há uma distância estrutural entre a proposta do
projeto e o sistema construído. Essa distância motiva a presente análise, que
tem por objetivos: (i) expor o referencial teórico das seis áreas de
conhecimento envolvidas; (ii) descrever o sistema implementado e seu pipeline
de classificação; (iii) apresentar o mapeamento projeto$\times$sistema e a
taxonomia de gaps; (iv) discutir convergências, limitações e ameaças à
validade; e (v) propor diretrizes de arquitetura para o alinhamento futuro.
A contribuição deste artigo é, portanto, a primeira caracterização
sistemática e fundamentada em evidência de código da lacuna entre uma proposta
de pesquisa de LLMs \textit{open source} para SDN e sua implementação
concreta.

% =====================================================================
\section{Referencial Teórico e Revisão da Literatura}
\label{sec:referencial}

A fundamentação deste trabalho articula seis eixos de conhecimento, sintetizados
a partir de revisão de literatura estruturada, com 122 referências efetivamente
citadas e validadas neste artigo.

\subsection{Método e Corpus de Análise}
\label{sec:metodo}

A seleção do corpus seguiu um procedimento de revisão integrativa
\citea{rA28}\citea{rA212}. Partiu-se de seis áreas de conhecimento
(Eixos~1--6), mapeadas a partir do Projeto de Pesquisa norteador
\citea{rRODRIGUES} e de relatórios de análise de código do sistema NOC
Lite: (1)~SDN/NFV; (2)~detecção e classificação de anomalias por ML
tradicional; (3)~LLMs \textit{open source} e modelos de linguagem para
NetOps; (4)~fine-tuning eficiente de LLMs (PEFT/LoRA/QLoRA); (5)~IA
explicável (XAI); e (6)~geração de cenários anômalos e \textit{datasets}
públicos. Como critérios de inclusão, admitiram-se obras primárias
(artigos, \textit{preprints} e relatórios técnicos) diretamente
relacionadas a um dos seis eixos e com relevância para a análise
arquitetural realizada; excluíram-se trabalhos fora do escopo de
detectação/mitigação de anomalias em redes e fontes sem validação
verificável. O corpus final compreende 122 fontes efetivamente citadas,
distribuídas pelos seis eixos, consolidadas em uma única lista de
referências ordenada alfabeticamente conforme a ABNT NBR 6023.

\subsection{Redes Definidas por Software e Virtualização de Funções de Rede}
\label{sec:sdn}

O paradigma SDN desacopla o plano de controle do plano de dados por meio de um
protocolo \textit{southbound} padronizado (OpenFlow) \citea{rA11} e de controladores
programáveis (NOX \citea{rA12}, POX, Ryu \citea{rA116}, OpenDaylight \citea{rA114}, ONOS \citea{rA115}).
Essa centralização lógica cria uma única fonte de verdade e uma superfície de
programação que, simultaneamente, facilita a automação e introduz novas
superfícies de ataque no plano de controle \citea{rA122}\citea{rA123}\citea{rA124}. A emulação
(Mininet \citea{rA14}, Open vSwitch \citea{rA15}) viabiliza a geração controlada de cenários
anômalos \citea{rA19}. Surveys definitivos consolidam a taxonomia e os desafios de
segurança da SDN \citea{rA19}\citea{rA110}.

Complementarmente, a Virtualização de Funções de Rede (NFV), institucionalizada
pelos \textit{white papers} e pelo \textit{framework} arquitetural da ETSI
\citea{rA127}\citea{rA128}\citea{rA129}, e revisada por Mijumbi et al. \citea{rA130}\citea{rA131} e Laghrissi \&
Taleb \citea{rA132}, estende o paradigma ao ciclo de vida de funções virtualizadas e
sua orquestração (MANO). Sobre esse alicerce, o \textit{Intent-Based
Networking} (IBN) emerge como a camada de gestão que traduz objetivos de alto
nível em configurações concretas e as garante por meio de laços fechados
\citea{rA133}. Trabalhos como Morales et al. \citea{rA134} e a survey de Panayiotou et al.
\citea{rA135} antecipam a reconfiguração proativa da topologia virtual orientada a
dados e ML. A detecção de anomalias em SDN transita, assim, de métodos
clássicos de ML \citea{rA126} para a integração de LLMs \textit{open source}
ajustados por LoRA/QLoRA.

\subsection{Detecção e Classificação de Anomalias por Métodos Tradicionais}
\label{sec:ml}

A literatura consolida que métodos estatísticos (cartas EWMA \citea{rA22}, Holt-Winters
\citea{rA21}, ARIMA/SARIMA \citea{rA25}\citea{rA27}) e classificadores de ML clássico (SVM
\citea{rA213}, \textit{Random Forest} \citea{rA210}, árvores de decisão \citea{rA217}, redes
rasas \citea{rA231}) formam o arcabouço de \textit{baselines} sobre o qual a detecção
moderna é avaliada. Surveys consagrados \citea{rA28}\citea{rA29}\citea{rA210}\citea{rA212} reforçam que
heurísticas de limiar e modelos estatísticos são insubstituíveis como
referência de comparação por serem interpretáveis e leves. Para os tipos de
anomalia do sistema, há corpus dedicado: DDoS é coberto por baselines de
árvore/RF \citea{rA222}\citea{rA224}\citea{rA225}; varredura de portas possui precedentes
heurísticos e de ML \citea{rA227}\citea{rA229}; exfiltração é abordada sobre metadados
de fluxo \citea{rA226}; e congestionamento/falha de enlace encontra raiz em detecção
estatística de séries temporais \citea{rA21}\citea{rA232}\citea{rA21}.

\subsection{LLMs Open Source e Modelos de Linguagem para NetOps}
\label{sec:llm}

Surveys dedicados mapeiam a aplicação de LLMs à operação e gestão de redes
(NO\&M/NetOps) \citea{rA31}\citea{rA32}\citea{rA33}\citea{rA34}\citea{rA35}\citea{rA36}\citea{rA37}. O ecossistema de modelos
\textit{open source} -- LLaMA/Llama 2 \citea{rA310}\citea{rA311}, Mistral 7B \citea{rA312}, Falcon
\citea{rA313}, Qwen/Qwen2.5 \citea{rA314}\citea{rA315}, BLOOM \citea{rA316} e, sobretudo, Gemma
\citea{rA317} (derivado da tecnologia Gemini) -- oferece pesos e licenças que tornam
viável a substituição de modelos fechados por alternativas locais e ajustáveis
\citea{rA41}. A arquitetura Transformer \citea{rA319}, o alinhamento por RLHF \citea{rA320} e o
\textit{Chain-of-Thought} (CoT) \citea{rA321} fundamentam, respectivamente, a
capacidade de raciocínio, de seguir instruções e de produzir justificativa
textual. Aplicações concretas incluem detecção de anomalias em SDN
\citea{rA325}\citea{rA327}, IBN orientado a LLM \citea{rA324}\citea{rA332}\citea{rA333} e agentes de rede
autônomos \citea{rA338}\citea{rA334}\citea{rA335}\citea{rA336}\citea{rA337}.

\subsection{Fine-Tuning Eficiente de LLMs (PEFT): LoRA, QLoRA}
\label{sec:peft}

O ajuste fino eficiente (PEFT) -- liderado por LoRA \citea{rA41} e estendido por QLoRA
\citea{rA42}, AdaLoRA \citea{rA43} e DoRA \citea{rA44} -- é a via consagrada para adaptar LLMs
\textit{open source} a domínios específicos com orçamento de hardware modesto,
reduzindo ordens de magnitude de parâmetros treináveis sem overhead de
inferência \citea{rA45}\citea{rA49}\citea{rA415}. Frameworks como LLM-Adapters \citea{rA411} e
LLaMA-Adapter \citea{rA412} demonstram que LLMs 7B ajustados por PEFT igualam ou
superam modelos 175B em \textit{zero-shot}. No domínio de tráfego de rede,
NetGPT \citea{rA425}, TrafficGPT \citea{rA426} e TrafficLLM \citea{rA427} mostram que LLMs podem ser
ajustados (com tokenização de domínio) para entender e classificar tráfego. A
base \textit{few-shot/zero-shot} (Brown \citea{rA422}, Wei \citea{rA321}, Achiam \citea{rA424})
legitima o uso de LLMs sem ajuste de pesos, mas a literatura documenta que
PEFT tipicamente supera o \textit{in-context learning} em custo e acurácia
\citea{rA49}.

\subsection{Inteligência Artificial Explicável (XAI)}
\label{sec:xai}

A XAI é requisito fundamental para a adoção confiável de detectores em domínios
de alta responsabilidade \citea{rA52}\citea{rA531}\citea{rA517}. Em NLP, o survey de Danilevsky et
al. \citea{rA51} sistematiza como modelos de linguagem produzem explicações; Ribeiro
et al. \citea{rA53} (LIME) e Lundberg \& Lee \citea{rA54} (SHAP) oferecem atribuição local.
O protótipo HuntGPT \citea{rA511} integra um detector de anomalias por ML a um LLM
que entrega justificativas textuais -- referência explícita do projeto. Estudos
recentes exploram LLMs para detecção de intrusão explicável \citea{rA520}\citea{rA521} e
explicabilidade em SDN \citea{rA522}. A calibração e a confiança das justificativas
são temas ativos \citea{rA531}\citea{rA532}\citea{rA533}\citea{rA534}.

\subsection{Geração de Cenários Anômalos e Datasets Públicos}
\label{sec:datasets}

A literatura estabelece UNSW-NB15 \citea{rA61}\citea{rA62} e CSE-CIC-IDS2018 \citea{rA63}\citea{rA64} como
benchmarks modernos, superando os obsoletos KDD99/NSL-KDD \citea{rA65}\citea{rA66}. A
geração dessas bases segue perfis comportamentais e emulação (Mininet,
GNS3, PNETLab) \citea{rA69}\citea{rA629}. Datasets SDN específicos (InSDN
\citea{rA610}, SNT, DDoS2019) e técnicas de \textit{data augmentation} (GAN/VAE/SMOTE)
endereçam o desbalanceamento de classes \citea{rA614}\citea{rA615}\citea{rA623}. A conversão de
atributos numéricos em representações textuais -- formalizada por AnoLLM
\citea{rA616} e revisada em surveys de LLMs em dados tabulares \citea{rA617}\citea{rA618} -- abre
caminho para detectores baseados em LLMs, eixo conectado diretamente à
inferência do sistema analisado \citea{rA630}.

% =====================================================================
\section{O Sistema Implementado: NOC Lite}
\label{sec:sistema}

O \textit{NOC Lite} é um sistema de monitoramento de disponibilidade e
performance de \textit{hosts} genéricos, estruturado como um monolito modular em
FastAPI com persistência em SQLite e um \textit{front-end} estático. A
arquitetura de contêineres (Figura~\ref{fig:arq}) cobre: \textit{front-end}
(HTML/JS); API backend (FastAPI/Uvicorn); motor de monitoramento
(\texttt{asyncio}); classificador de anomalias; simulador; notificador
(Telegram); e banco de dados SQLite. Atenção especial: o classificador de
anomalias \textbf{não participa do laço de monitoramento} -- é acionado sob
demanda via \texttt{GET /api/anomalies/classify}, jamais pelo
\textit{MonitorEngine}.

\begin{figure}[H]
  \centering
  \includegraphics[width=0.85\textwidth]{fig_arquitetura_sistema}
  \caption{Arquitetura de contêineres do NOC Lite (C4, nível contêiner).
  Fonte: elaborada a partir de \texttt{arquitetura\_sistema.md} (Dan).}
  \label{fig:arq}
\end{figure}

\subsection{Pipeline de Classificação}
\label{sec:pipeline}

O subsistema de classificação expõe três estratégias sobre um conjunto fechado
de seis classes -- \texttt{NORMAL}, \texttt{DDOS\_ATTACK}, \texttt{PORT\_SCAN},
\texttt{HOST\_DOWN}, \texttt{DATA\_EXFILTRATION} e \texttt{NETWORK\_CONGESTION}
-- conforme o fluxo da Figura~\ref{fig:pipe}.

\begin{figure}[H]
  \centering
  \includegraphics[width=0.78\textwidth]{fig_pipeline_classificacao}
  \caption{Pipeline de classificação de anomalias do NOC Lite: extração de
  features (janela de 5~min), três classificadores paralelos e tokenização
  para o LLM. Fonte: \texttt{notas\_coder.md} (John).}
  \label{fig:pipe}
\end{figure}

\begin{enumerate}
  \item \textbf{Heurístico (baseline 1)} -- \texttt{classify\_heuristic}:
  regras de limiar estático. Ex.: \texttt{network\_in\_bps > 50\,Mb/s} e
  \texttt{cpu\_usage > 75} $\rightarrow$ \texttt{DDOS\_ATTACK};
  \texttt{ping\_latency > 120\,ms} $\rightarrow$ \texttt{NETWORK\_CONGESTION}.
  Produz confiança fixa e explicação textual.
  \item \textbf{Machine Learning (baseline 2)} -- \texttt{classify\_ml}:
  \texttt{DecisionTreeClassifier(max\_depth=5)} do scikit-learn, treinado
  \textbf{exclusivamente com dados sintéticos} (1.200 amostras, 200 por
  classe). Sem PEFT, sem LLM.
  \item \textbf{IA Generativa (Gemini)} -- \texttt{classify\_generative\_ai}:
  monta \textit{prompt} estruturado com \textit{snapshot} tokenizado e envia ao
  modelo fechado \texttt{gemini-1.5-flash}, retornando JSON com
  \texttt{classification}, \texttt{confidence} e \texttt{explanation}. Possui
  \textit{fallback} para a heurística.
\end{enumerate}

A função \texttt{tokenise\_metrics} converte o vetor numérico de 10 \textit{
features} em texto estruturado (ex.: \texttt{PingSuccessRate: 0.95 |
AvgPingLatency: 12.3 ms | CPUUsage: 34.1\%}), espelhando a recomendação do
projeto de converter atributos numéricos em representações textuais
\citea{rA616}\citea{rA431}\citea{rA432}. Esse texto é a ponte que permite ao LLM ``ler'' os limiares
de rede em linguagem natural.

\subsection{Simulador}
\label{sec:sim}

O \texttt{anomaly\_simulator.py} gera cenários anômalos sintéticos
(\texttt{generate\_simulated\_metrics}) para as seis classes via
\texttt{random.uniform}, e injeta esses vetores diretamente no banco
(\texttt{run\_simulated\_check\_cycle}), podendo popular histórico retroativo.
O agendamento em tempo real usa estado em memória. Trata-se de dados 100\%
sintéticos, sem dependência de datasets públicos.

% =====================================================================
\section{Mapeamento Projeto $\times$ Sistema e Gaps}
\label{sec:mapeamento}

A Tabela~\ref{tab:mapa} sintetiza a correspondência entre os objetivos/componentes
do projeto norteador \citea{rRODRIGUES} e o que o NOC Lite efetivamente entrega, com
o respectivo status de alinhamento. A Figura~\ref{fig:mapa} ilustra essa
relação.

\begin{longtable}{cp{5.4cm}p{5.0cm}c}
\caption{Mapeamento Projeto de Pesquisa $\times$ NOC Lite.}
\label{tab:mapa}\\
\toprule
\textbf{\#} & \textbf{Objetivo/Componente do Projeto} & \textbf{Entrega do NOC Lite} & \textbf{Status} \\
\midrule
\endfirsthead
\toprule
\textbf{\#} & \textbf{Objetivo/Componente do Projeto} & \textbf{Entrega do NOC Lite} & \textbf{Status} \\
\midrule
\endhead
1 & LLM \textit{open source} (Llama/Falcon) com fine-tuning & Gemini 1.5 Flash (fechado); Árvore de Decisão & Divergente \\
2 & Fine-tuning PEFT (LoRA/QLoRA) & Inexistente; ML é árvore treinada do zero & Ausente \\
3 & Ambiente SDN virtualizado (Mininet/OVS) & Hosts IP genéricos; sondas ICMP/TCP/HTTP/DNS/SNMP & Ausente \\
4 & Mitigação via reconfiguração VNT/roteamento (ciclo fechado) & Apenas detecção + alerta Telegram; sem ação corretiva & Ausente \\
5 & Anomalias de performance (congestionamento, falha de enlace) & Mistura anomalias de segurança/operacionais com congest. genérico & Parcial \\
6 & Datasets públicos (UNSW-NB15, CIC-IDS2018, Synthetic SDN) & Dados 100\% sintéticos (random.uniform) & Divergente \\
7 & Baselines de comparação (heurística / ML) & Heurística + DecisionTree & Convergente \\
8 & Tokenização de métricas em texto estruturado & \texttt{tokenise\_metrics} & Convergente \\
9 & Justificativa textual (XAI) & \texttt{explanation} (com modelo fechado) & Convergente \\
10 & Avaliação de tempo de resposta por classificador & \texttt{latency\_ms} medido por classificador & Convergente \\
11 & Métricas resumidas (throughput, latência, jitter, perda) & SLA/jitter/tendência/health-score + SNMP & Convergente \\
12 & Ciclo fechado coleta$\rightarrow$classificação$\rightarrow$mitigação & Loop aberto: coleta$\rightarrow$severidade$\rightarrow$alerta & Aberto \\
\bottomrule
\end{longtable}

\begin{figure}[H]
  \centering
  \includegraphics[width=0.85\textwidth]{fig_mapeamento_artigo_sistema}
  \caption{Mapeamento visual entre os objetivos do artigo norteador e o
  sistema NOC Lite. Fonte: elaborada por Chris.}
  \label{fig:mapa}
\end{figure}

\subsection{Taxonomia de Gaps}
\label{sec:gaps}

A análise arquitetural identifica sete gaps, seis de severidade crítica/maior
(Figura~\ref{fig:gaps}). Cada gap é sustentado por evidência de código.

\begin{longtable}{cp{4.0cm}p{8.0cm}c}
\caption{Taxonomia de gaps (NOC Lite $\times$ Projeto).}
\label{tab:gaps}\\
\toprule
\textbf{ID} & \textbf{Gap} & \textbf{Evidência no código / Descrição} & \textbf{Severidade} \\
\midrule
\endfirsthead
\toprule
\textbf{ID} & \textbf{Gap} & \textbf{Evidência no código / Descrição} & \textbf{Severidade} \\
\midrule
\endhead
GAP-1 & SDN vs. hosts genéricos & Busca por \texttt{sdn|mininet|vnt|controller} $\rightarrow$ 0 ocorrências; \texttt{Host.address}arbítrio & Crítico \\
GAP-2 & Gemini fechado vs. LLM \textit{open source} & \texttt{genai.GenerativeModel("gemini-1.5-flash")}; depende de \texttt{GEMINI\_API\_KEY} & Crítico \\
GAP-3 & Ausência de LoRA/QLoRA & Busca por \texttt{lora|qlora|peft|fine.?tun} $\rightarrow$ 0; único modelo treinável é DecisionTree & Crítico \\
GAP-4 & Ausência de mitigação / loop aberto & \texttt{scheduler.py} não referencia nenhum \texttt{classify\_*}; inferência só em \texttt{/api/anomalies/classify} & Crítico \\
GAP-5 & Dados sintéticos vs. datasets públicos & \texttt{random.uniform} em \texttt{anomaly\_simulator.py}; busca \texttt{unsw|cic-ids} $\rightarrow$ 0 & Maior \\
GAP-6 & Performance vs. segurança/operacionais & Classes majoritariamente de segurança; congest. por limiar genérico de latência de ping & Maior \\
GAP-7 & Interpretabilidade com modelo fechado & \texttt{explanation} provém do Gemini; sem rastreabilidade de pesos nem reprodutibilidade local & Menor \\
\bottomrule
\end{longtable}

\begin{figure}[H]
  \centering
  \includegraphics[width=0.82\textwidth]{fig_gaps}
  \caption{Taxonomia de gaps estruturais entre o projeto norteador e o NOC
  Lite. Fonte: elaborada por Chris.}
  \label{fig:gaps}
\end{figure}

\textbf{GAP-1 (SDN vs. hosts).} O projeto fundamenta toda a proposta no
paradigma SDN; o NOC Lite monitora hosts IP/DNS arbitrários, sem controlador,
topologia virtualizada ou planos separados \citea{rA19}. \textbf{GAP-2 (modelo
fechado).} O projeto exige LLMs \textit{open source} e classifica explicitamente
o Gemini como modelo fechado \citea{rRODRIGUES}; o sistema o usa em regime \textit{
zero-shot} \citea{rA317}. \textbf{GAP-3 (sem LoRA/QLoRA).} A seção de PEFT do projeto
não tem correspondente: o único modelo treinável é uma árvore de decisão do
zero \citea{rA41}\citea{rA42}. \textbf{GAP-4 (loop aberto).} A classificação por IA não está
integrada ao laço de detecção; o ciclo observa e alerta, sem atuação autônoma
\citea{rA133}. \textbf{GAP-5 (dados sintéticos).} O treino e o \textit{few-shot}
repousam no simulador; ausente a validação cruzada com UNSW-NB15,
\citea{rA512}. \textbf{GAP-6 (escopo).} O foco do
projeto é performance (congestionamento, falha de enlace, degradação de
desempenho); o sistema privilegia classes de segurança \citea{rA512}. \textbf{GAP-7
(XAI com fechado).} A justificativa provém de modelo opaco, fragilizando a
promessa de XAI do projeto \citea{rA511}\citea{rA52}.

% =====================================================================
\section{Discussão}
\label{sec:discussao}

\subsection{Convergências preservadas}
\label{sec:convergencias}

Apesar dos desalinhamentos, o NOC Lite honra parcela da intenção do projeto:

\begin{enumerate}
  \item \textbf{Baselines de comparação} -- implementa exatamente heurística e
  árvore de decisão, permitindo comparação direta lado a lado \citea{rA210}\citea{rA217}.
  \item \textbf{Tokenização de métricas} -- \texttt{tokenise\_metrics} converte
  o vetor numérico em representação semiestruturada, alinhada à recomendação de
  conversão textual \citea{rA616}\citea{rA431}.
  \item \textbf{Justificativa textual / XAI} -- todos os classificadores retornam
  \texttt{explanation}, atendendo ao objetivo qualitativo de interpretabilidade
  \citea{rA511}\citea{rA537}.
  \item \textbf{Métricas de tempo de resposta} -- o endpoint mede
  \texttt{latency\_ms} de cada classificador, entregando a métrica de desempenho
  exigida \citea{rA31}.
  \item \textbf{Métricas resumidas de rede} -- SLA, jitter, tendência e
  health-score capturam latência/perda/jitter, coincidentes com as métricas de
  performance do projeto \citea{rA19}.
\end{enumerate}

\subsection{Limitações e ameaças à validade}
\label{sec:ameacas}

\textbf{Viés de dados sintéticos.} O classificador de ML é treinado 100\% com o
próprio simulador, que compartilha os mesmos limiares da heurística; a árvore
aprende essencialmente a ``copiar'' a heurística em forma de árvore, com alto
desempenho no sintético e provável colapso em produção \citea{rA414}\citea{rA615}.
\textbf{Confiança não calibrada.} A confiança do Gemini é autorreportada e a do
ML é probabilidade de folha não calibrada \citea{rA532}; a comparação entre
classificadores é, portanto, metodologicamente assimétrica. \textbf{
Orquestração} ausente: a classificação está fora do ciclo de monitoramento,
impossibilitando detecção contínua. \textbf{Validade externa}: a ausência de
datasets públicos impede a comparação com o estado da arte \citea{rA65}\citea{rA630}.

\textbf{Ameaças à validade interna} incluem a dependência de API externa paga
(Gemini) e a fragilidade do \textit{parser} JSON (apenas checagem de presença
de chave). \textbf{Ameaças à validade de construção}: os limiares de 50~Mb/s e
120~ms são fixos e universais, sem \textit{baseline} por host/link, gerando
falsos positivos em enlaces de capacidades distintas \citea{rA21}.

% =====================================================================
\section{Diretrizes de Arquitetura para Alinhamento}
\label{sec:diretrizes}

Para reduzir a distância entre o NOC Lite e o projeto, propõe-se uma evolução
arquitetural incremental, respeitando o que já existe e fechando os sete gaps.
O ponto central é a transição de um \textit{loop aberto} para um \textit{ciclo
fechado} com LLM \textit{open source}, LoRA e mitigação, conforme a
Figura~\ref{fig:ciclo}.

\begin{figure}[H]
  \centering
  \includegraphics[width=0.85\textwidth]{fig_ciclo_fechado_proposto}
  \caption{Arquitetura proposta de ciclo fechado: coleta (SDN + hosts) $\rightarrow$
  tokenização $\rightarrow$ classificação (LLM \textit{open source} + LoRA)
  $\rightarrow$ mitigação (módulo de ação) $\rightarrow$ reconfiguração via
  controlador. Fonte: elaborada por Chris.}
  \label{fig:ciclo}
\end{figure}

\begin{enumerate}
  \item \textbf{Substituição do modelo fechado por LLM local (fecha GAP-2,
  GAP-3).} Introduzir contêiner de inferência local (vLLM/llama.cpp/Ollama)
  servindo Llama 3 (7--13B) ou Falcon via endpoint compatível com OpenAI;
  substituir \texttt{gemini-1.5-flash} por cliente HTTP local, mantendo o
  \textit{fallback} heurístico. Aplicar LoRA/QLoRA (rank 16--32, alpha 32)
  sobre o modelo base, reutilizando os \textit{prompts} de \texttt{
  tokenise\_metrics} \citea{rA41}\citea{rA42}\citea{rA312}.
  \item \textbf{Pipeline de dados com datasets públicos (fecha GAP-5).} Criar
  módulo \texttt{data\_ingest} que baixa/normaliza UNSW-NB15, CSE-CIC-IDS2018
  e Synthetic SDN para o \textit{schema} de \textit{features}; treinar LoRA e a
  árvore com dados reais, preservando o simulador para demonstração \citea{rA61}\citea{rA64}.
  \item \textbf{Ambiente SDN e fechamento do ciclo (fecha GAP-1, GAP-4).}
  Adicionar contêiner Mininet + controlador Ryu que emule a topologia e exponha
  telemetria de enlace/fluxo; generalizar \texttt{Host}/\texttt{CheckResult}
  para telemetria de fluxo SDN; criar \texttt{mitigation\_engine} que mapeia
  \texttt{classification} $\rightarrow$ ação corretiva (re-rota, QoS,
  \textit{drop} de fluxo) via NorthBound API do controlador; integrar o
  classificador dentro do \texttt{MonitorEngine}, fechando o ciclo \citea{rA14}\citea{rA116}\citea{rA133}.
  \item \textbf{Realinhamento do catálogo de anomalias (fecha GAP-6).}
  Separar claramente performance SDN (congestionamento de enlace, falha de
  enlace, degradação de throughput/jitter) de segurança/operacional (DDoS, port
  scan, exfiltração, host down), priorizando a primeira família \citea{rA512}.
  \item \textbf{Preservação dos pontos fortes.} Manter \texttt{tokenise\_metrics},
  o par baseline heurístico + ML, a medição de \texttt{latency\_ms} e o formato
  de \texttt{explanation} como base do novo pipeline \citea{rA616}\citea{rA431}.
  \item \textbf{Roadmap de alinhamento} (espelhando o cronograma do projeto):
  (1) contêiner de LLM \textit{open source} local + \texttt{tokenise\_metrics}
  $\rightarrow$ \textit{zero-shot}; (2) ingestão dos 3 datasets + treino
  LoRA/QLoRA; (3) Mininet + controlador + telemetria de fluxo; (4)
  \texttt{mitigation\_engine} + integração da classificação no
  \texttt{MonitorEngine}; (5) fechamento do ciclo e avaliação comparativa
  (heurística vs. ML vs. LLM \textit{open source} com/sem LoRA), reutilizando
  \texttt{latency\_ms} e acrescentando acurácia/consistência.
\end{enumerate}

% =====================================================================
\section{Conclusão}
\label{sec:conclusao}

O presente artigo correlacionou o projeto de pesquisa que propõe detecção e
mitigação de anomalias em SDN com LLMs \textit{open source} ajustados por
PEFT/LoRA \citea{rRODRIGUES} com o sistema efetivamente implementado, o NOC Lite.
Evidenciou-se que o sistema realiza com sucesso a parcela de \textit{
baselines} e de representação textual do projeto -- heurística, árvore de
decisão, tokenização de métricas, justificativa textual e medição de latência
-- mas diverge estruturalmente nos eixos centrais da proposta. Especificamente,
o NOC Lite (i) monitora hosts genéricos e não SDN; (ii) emprega o modelo
fechado Gemini em vez de LLM \textit{open source}; (iii) não aplica LoRA/QLoRA;
(iv) não realiza mitigação nem opera em ciclo fechado; e (v) depende
exclusivamente de dados sintéticos, sem validação cruzada com datasets
públicos.

Tais gaps não são ajustes pontuais: exigem evolução arquitetural substantiva,
sintetizada nas diretrizes de um ciclo fechado com LLM local, LoRA e módulo de
mitigação. As convergências identificadas funcionam como alicerce seguro sobre
o qual construir essa evolução, uma vez que o \texttt{tokenise\_metrics}, os
baselines e a medição de latência já cumprem objetivos específicos do projeto.
Como trabalho futuro, recomenda-se a implementação do \texttt{mitigation\_engine}
e a condução do experimento comparativo head-to-head entre LLM \textit{open
source} ajustado por QLoRA e LLM fechado \textit{zero-shot} sobre o mesmo
conjunto de métricas de rede -- lacuna exatamente apontada pela literatura
\citea{rA49}\citea{rA630}.

% =====================================================================
% Referencias (ABNT manual, geradas a partir das areas de literatura)
\section*{Referências}
% restored numeric
\input{referencias.tex}


% =====================================================================
% Referências ABNT (autor-data)
% =====================================================================
\begin{thebibliography}{99}
\raggedright
\bibitem{rA11} McKEOWN et al, 2008. A network in a laptop: rapid prototyping for software-defined networks.
\bibitem{rA110} XIA et al, 2015. Chain-of-Thought Prompting Elicits Reasoning in Large Language Models.
\bibitem{rA114} MEDVED et al, 2014. MEDVED, J.; VARGA, R.; TKACIK, A.; GRAY, K. OpenDaylight: Towards a Model-Driven SDN Controller Architecture. ACM SIGCOMM HotSDN '14, p. 1–6, 2014..
\bibitem{rA115} BERDE et al, 2014. BERDE, P.; GEROLA, M.; HART, J.; HIGUCHI, Y.; KOBAYASHI, M.; KOIDE, T.; et al. ONOS: Towards an Open, Distributed SDN OS. ACM SIGCOMM HotSDN '14, p. 1–6, 2014..
\bibitem{rA116} NTT Laboratories. Ryu SDN Framework - Open-source SDN Platform Software. NTT Technical Review, 2014. NTT Laboratories. Ryu SDN Framework — Open-source SDN Platform Software. NTT Technical Review, v. 12, n. 8, 2014..
\bibitem{rA12} GUDE et al, 2008. GUDE, N.; KOPONEN, T.; PETTIT, J.; PFAFF, B.; MCKEOWN, N.; SHENKER, S.; CASADO, M. NOX: Towards an Operating System for Networks. ACM SIGCOMM Computer Communication Review, v. 38, n. 3, p. 105–110, 2008..
\bibitem{rA122} PORRAS et al, 2012. PORRAS, P.; SHIN, S.; YEGNESWARAN, V.; FONG, M.; TYSON, M.; GU, G. A Security Enforcement Kernel for OpenFlow Networks. ISOC NDSS '12, 2012..
\bibitem{rA123} LI et al, 2016. Referencia: LI et al. (2016).
\bibitem{rA124} BHUIYAN et al, 2023. BHUIYAN, Z. A.; ISLAM, S.; ISLAM, M. M.; ULLAH, A. A.; NAZ, F.; RAHMAN, M. S. On the (in)Security of the Control Plane of SDN Architecture: A Survey. IEEE Access, v. 11, p. 91550–91582, 2023..
\bibitem{rA126} LATIF et al, 2022. LATIF, Z.; SHARIF, K.; LI, F.; KARIM, M. M.; BELGHITI, M. A.; SANGEEN, A. A Machine Learning-Based Anomaly Prediction Service for Software-Defined Networks. Sensors, v. 22, n. 21, artigo 8434, 2022..
\bibitem{rA127} CHIOSI et al, 2012. CHIOSI, M.; CLARKE, D.; HAMPSON, P.; et al.
\bibitem{rA128} ETSI NFV ISG. GS NFV 002 V1.1.1: Network Functions Virtualisation (NFV); Architectural Framework. ETSI et al, 2013. ETSI, arquitetura, MANO, orquestração, VNF placement. E5. Intent-Based Networking (IBN), reconfiguração proativa de topologia virtual (VNT) e roteamento em ciclo fechado (closed-loop.
\bibitem{rA129} ETSI NFV ISG. GS NFV-MAN 001 V1.1.1: Network Functions Virtualisation (NFV); Management and Orchestration. ETSI et al, 2014. ETSI, arquitetura, MANO, orquestração, VNF placement. E5. Intent-Based Networking (IBN), reconfiguração proativa de topologia virtual (VNT) e roteamento em ciclo fechado (closed-loop.
\bibitem{rA130} MIJUMBI et al, 2015. MIJUMBI, R.; SERRAT, J.; GORRICHO, J.-L.; BOUTEN, N.; DE TURCK, F.; BOUTABA, R. Network Function Virtualization: State-of-the-Art and Research Challenges. IEEE Communications Magazine, v. 53, n. 2, p. 216–228, 2015..
\bibitem{rA131} MIJUMBI et al, 2016. MIJUMBI, R.; SERRAT, J.; GORRICHO, J.-L.; BOUTEN, N.; DE TURCK, F.; BOUTABA, R. Network Function Virtualization: State-of-the-Art and Research Challenges. IEEE Communications Magazine, v. 53, n. 2, p. 216–228, 2015..
\bibitem{rA132} LAGHRISSI et al, 2019. LAGHRISSI, A.; TALEB, T. A Survey on the Placement of Virtual Resources and Virtual Network Functions. IEEE Communications Surveys & Tutorials, v. 21, n. 2, p. 1409–1434, 2019..
\bibitem{rA133} LEIVADEAS et al, 2022. LEIVADEAS, A.; FALKNER, M. A Survey on Intent-Based Networking. IEEE Communications Surveys & Tutorials, v. 25, n. 1, p. 625–655, 1º trim. 2023.
\bibitem{rA134} MORALES et al, 2017. MORALES, F.; RUIZ, M.; CONTRERAS, L. M.; FERRARIS, G.; STAVDAS, A.; CUGINI, F.; VELASCO, L. Virtual Network Topology Adaptability Based on Data Analytics for Traffic Prediction. IEEE/OSA Journal of Optical Communications and Networking, v. 9, n. 1, p. A35–A45, 2017..
\bibitem{rA135} VELASCO et al, 2022. roteamento em ciclo fechado.
\bibitem{rA14} LANTZ et al, 2010. A network in a laptop: rapid prototyping for software-defined networks.
\bibitem{rA15} PFAFF et al, 2015. PFAFF, B.; MCKEOWN, N.; SHENKER, S.; CASADO, M. NOX: Towards an Operating System for Networks. ACM SIGCOMM Computer Communication Review, v. 38, n. 3, p. 105–110, 2008..
\bibitem{rA19} KREUTZ et al, 2015. KREUTZ, D.; RAMOS, F. M. V.; VERISSIMO, P. E.; ROTHENBERG, C. E.; AZODOLMOLKY, S.; UHLIG, S. Software-Defined Networking: A Comprehensive Survey. Proceedings of the IEEE, v. 103, n. 1, p. 14–76, 2015..
\bibitem{rA21} BRUTLAG, 2000. [Acesso via assinatura].
\bibitem{rA210} BHUYAN et al, 2014. BHUYAN, M. H.; BHATTACHARYYA, D. K.; KALITA, J. K. Network Anomaly Detection: Methods, Systems and Tools. IEEE Communications Surveys & Tutorials, v. 16, n. 1, p. 303–336, 2014. DOI: https://doi.org/10.1109/SURV.2013.052213.00046 URL: https://ieeexplore.ieee.org/iel7/9739/6734841/06524462.pdf.
\bibitem{rA212} BUCZAK et al, 2016. BUCZAK, A. L.; GÜVEN, E. A Survey of Data Mining and Machine Learning Methods for Cyber Security Intrusion Detection. IEEE Communications Surveys & Tutorials, v. 18, n. 2, p. 1153–1176, 2016. DOI: https://doi.org/10.1109/COMST.2015.2494502 (>5300 citações) URL: https://ieeexplore.ieee.org/document/7307098/.
\bibitem{rA213} CORTES et al, 1995. CORTES, C.; VAPNIK, V. Support-Vector Networks. Machine Learning, v. 20, n. 3, p. 273–297, 1995. DOI: https://doi.org/10.1007/BF00994018 URL: https://link.springer.com/article/10.1007/BF00994018.
\bibitem{rA217} CARUANA et al, 2006. CARUANA, R.; NICULESCU-MIZIL, A. An Empirical Comparison of Supervised Learning Algorithms. In: Proc. of the 23rd International Conference on Machine Learning (ICML 2006), p. 161–168, 2006. DOI: https://doi.org/10.1145/1143844.1143865 URL: https://www.cs.cornell.edu/~caruana/ctp/ct.papers/caruana.icml06.pdf.
\bibitem{rA22} ROBERTS, 1959. ROBERTS, S. W. Control Chart Tests Based on Geometric Moving Averages. Technometrics, v. 1, n. 3, p. 239–250, 1959..
\bibitem{rA222} HODO et al, 2017. HODO, E.; BELLEKENS, X.; HAMILTON, A.; DUBOULIL, C.; PHAM, Q. N. Threat Analysis of IoT Networks Using Artificial Neural Network Intrusion Detection System. In: Int. Symp. on Networks, Computers and Communications (ISNCC), IEEE, 2017. DOI: https://doi.org/10.1109/ISNCC.2017.8005728 URL: https://ieeexplore.ieee.org/document/8005728/.
\bibitem{rA224} ALDUAILIJ et al, 2022. [Acesso via assinatura].
\bibitem{rA225} NEGANDHI et al, 2016. NEGANDHI, P.; TRIVEDI, Y. Intrusion Detection System Using Random Forest on the NSL-KDD Dataset. In: Proc. of Int. Conf. on Computing, Analytics and Security Trends (CAST), 2016. URL: https://www.researchgate.net/publication/335718695_Intrusion_Detection_System_Using_Random_Forest_on_the_NSL-KDD_Dataset.
\bibitem{rA226} WILLEMS et al, 2023. [Acesso via assinatura].
\bibitem{rA227} SHEHAB et al, 2024. SHEHAB, R. et al. Improving Port Scan Cybersecurity Risks Detection using Feature Selection and Machine Learning. JATIT, v. 102, n. 16, 2024. URL: https://www.jatit.org/volumes/Vol102No16/9Vol102No16.pdf.
\bibitem{rA229} STANIFORD-CHEN et al, 1996. STANIFORD-CHEN, S.; CHEUNG, S.; CRAWFORD, R.; DILGER, M.; FRANK, J.; HOAGLAND, J.; LEVITT, K.; WEINTRAUB, C.; YIP, D.; ZERKLE, D. GrIDS — A Graph-Based Intrusion Detection System for Large Networks. In: Proc. of the 19th National Information Systems Security Conference (NISSC), 1996. URL: https://seclab.cs.ucdavis.edu/papers/nissc96.pdf.
\bibitem{rA231} CANNADY, 1998. CANNADY, J. D. Artificial Neural Networks for Misuse Detection. In: Proc. of the 1998 National Information Systems Security Conference (NISSC), 1998. URL: https://csrc.nist.gov/nissc/1998/proceedings/paperF13.pdf.
\bibitem{rA232} LAKHINA et al, 2004. LAKHINA, A.; CROVELLA, M.; DIOT, C. Diagnosing Network-Wide Traffic Anomalies. In: Proc. of ACM SIGCOMM 2004, p. 219–230, 2004. DOI: https://doi.org/10.1145/1030194.1015492 URL: https://dl.acm.org/doi/10.1145/1030194.1015492.
\bibitem{rA25} MOAYEDI et al, 2008. MOAYEDI, H. Z.; AMINI, M. ARIMA Model for Network Traffic Prediction and Anomaly Detection. In: Proc. of the 4th International Conference on Information Technology (ICIT 2008), IEEE, 2008. DOI: https://doi.org/10.1109/ICIT.2008.22 (via IEEE Xplore) URL: https://ieeexplore.ieee.org/document/4631947/.
\bibitem{rA27} FONDO-FERREIRO et al, 2026. FONDO-FERREIRO, P. et al. Real-time SARIMA-based anomaly detection in operator networks. Computer Networks (Elsevier), 2026. DOI: [NÃO VERIFICADO — usar URL abaixo; artigo de 2026 ainda em indexação] URL: https://www.sciencedirect.com/science/article/pii/S1389128626000824.
\bibitem{rA28} CHANDOLA et al, 2009. CHANDOLA, V.; BANERJEE, A.; KUMAR, V. Anomaly Detection: A Survey. ACM Computing Surveys, v. 41, n. 3, p. 1–58, 2009. DOI: https://doi.org/10.1145/1541880.1541882 URL: https://dl.acm.org/doi/10.1145/1541880.1541882.
\bibitem{rA29} AHMED et al, 2016. A Survey on Large Language Models for Communication, Network, and Service Management: Application Insights, Challenges, and Future Directions.
\bibitem{rA31} LIU et al, 2025. On IoT intrusion detection based on data augmentation for enhancing learning on unbalanced samples.
\bibitem{rA310} TOUVRON et al, 2023. TOUVRON, H.; LAVRIL, T.; IZACARD, G.; MARTINET, X.; LACHAUX, M.-A.; LACROIX, T. et al. LLaMA: Open and Efficient Foundation Language Models. 2023. arXiv:2302.13971. DOI: https://doi.org/10.48550/arXiv.2302.13971 URL: https://arxiv.org/abs/2302.13971.
\bibitem{rA311} TOUVRON et al, 2023. TOUVRON, H.; LAVRIL, T.; IZACARD, G.; MARTINET, X.; LACHAUX, M.-A.; LACROIX, T. et al. LLaMA: Open and Efficient Foundation Language Models. 2023. arXiv:2302.13971. DOI: https://doi.org/10.48550/arXiv.2302.13971 URL: https://arxiv.org/abs/2302.13971.
\bibitem{rA312} JIANG et al, 2023. JIANG, A. Q.; SABLAYROLLES, A.; MENSCH, A.; BAMFORD, C.; CHAPLOT, D. S. et al. Mistral 7B. 2023. arXiv:2310.06825. DOI: https://doi.org/10.48550/arXiv.2310.06825 URL: https://arxiv.org/abs/2310.06825.
\bibitem{rA313} ALMAZROUEI et al, 2023. The Falcon Series of Open Language Models.
\bibitem{rA314} BAI et al, 2023. Qwen Technical Report.
\bibitem{rA315} YANG et al, 2024. YANG, Z.; TANG, J. P-Tuning v2: Prompt Tuning Can Be Comparable to Fine-tuning Universally Across Scales and Tasks. ACL, 2022. arXiv:2110.07602. DOI: https://doi.org/10.48550/arXiv.2110.07602 URL: https://arxiv.org/abs/2110.07602.
\bibitem{rA316} BIGSCIENCE WORKSHOP et al, 2022. BLOOM: A 176B-Parameter Open-Access Multilingual Language Model.
\bibitem{rA317} GEMMA TEAM (Google DeepMind). Gemma: Open Models Based on Gemini Research and Technology. arXiv preprint arXiv:2403.08295, 2024. Gemma: Open Models Based on Gemini Research and Technology.
\bibitem{rA319} VASWANI et al, 2017. Attention Is All You Need.
\bibitem{rA32} LIU et al, 2024. On IoT intrusion detection based on data augmentation for enhancing learning on unbalanced samples.
\bibitem{rA320} OUYANG et al, 2022. Training Language Models to Follow Instructions with Human Feedback (InstructGPT).
\bibitem{rA321} WEI et al, 2022. Chain-of-Thought Prompting Elicits Reasoning in Large Language Models.
\bibitem{rA324} HOSSAIN et al, 2025. NetIntent: Leveraging Large Language Models for End-to-End Intent-Based SDN Automation.
\bibitem{rA325} SDN-Xplain: LLM-Driven Anomaly Detection and Human-Interpretable Explanation in SDN. SKKU IoT Lab / KICS Winter 2026 (preprint), 2026. Referencia: SDN-Xplain: LLM-Driven Anomaly Detection and Human-Interpretable Explanation in SDN. SKKU IoT Lab / KICS Winter 2026 (preprint). (2026).
\bibitem{rA327} Lightweight Fine-Tuning of LLMs for Explainable Intrusion Detection in SDN. arXiv/ResearchGate preprint, 2025. Lightweight Fine-Tuning of LLMs for Explainable Intrusion Detection in SDN." arXiv:2510.23313 (preprint), 2025. [título/autoria consolidados no preprint] URL: https://arxiv.org/html/2510.23313v1.
\bibitem{rA33} BOATENG et al, 2024. A Survey on Large Language Models for Communication, Network, and Service Management: Application Insights, Challenges, and Future Directions.
\bibitem{rA332} ADANZA et al, 2025. Leveraging Generative AI for Intent-Based Networking.
\bibitem{rA333} Towards Intent-Based Network Management: Large Language Models for Intent Extraction in 5G Core Networks. IEEE / GenAI Net research library, 2024. Towards Explainable Network Intrusion Detection using Large Language Models." arXiv:2408.04342, 2024. URL: https://arxiv.org/abs/2408.04342.
\bibitem{rA334} Exploring LLM-Based Agents for Root Cause Analysis. ACM SIGMETRICS / arXiv:2403.04123, 2024. Exploring LLM-Based Agents for Root Cause Analysis". ACM SIGMETRICS / arXiv:2403.04123, 2024..
\bibitem{rA335} Empowering AIOps: Leveraging Large Language Models for IT Operations Management. arXiv:2501.12461, 2025. Empowering AIOps: Leveraging Large Language Models for IT Operations Management". arXiv:2501.12461, 2025..
\bibitem{rA336} WANG et al, 2025. On the Calibration of Large Language Models and Alignment.
\bibitem{rA337} LLM-Driven Multi-Agent Framework for Autonomous Network Slice Management. IEEE, 2026. Large Language Models Meet Next-Generation Networking Technologies: A Review.
\bibitem{rA338} CHEN et al, 2024. Explainable Intrusion Detection Systems (X-IDS): A Survey of Current Methods, Challenges, and Opportunities.
\bibitem{rA34} HANG et al, 2024. Large Language Models Meet Next-Generation Networking Technologies: A Review.
\bibitem{rA35} HUANG et al, 2024. Explainable and Interpretable Multimodal Large Language Models: A Comprehensive Survey.
\bibitem{rA36} HONG et al, 2025. A Comprehensive Survey on LLM-Based Network Management and Orchestration.
\bibitem{rA37} ZHANG et al, 2025. When LLMs Meet Cybersecurity: A Systematic Literature Review.
\bibitem{rA39} TANG et al, 2024. TANG, J. P-Tuning v2: Prompt Tuning Can Be Comparable to Fine-tuning Universally Across Scales and Tasks. ACL, 2022. arXiv:2110.07602. DOI: https://doi.org/10.48550/arXiv.2110.07602 URL: https://arxiv.org/abs/2110.07602.
\bibitem{rA41} HU et al, 2021. Referencia: HU et al. (2021).
\bibitem{rA411} HU et al, 2023. Referencia: HU et al. (2023).
\bibitem{rA412} ZHANG et al, 2024. When LLMs Meet Cybersecurity: A Systematic Literature Review.
\bibitem{rA414} DING et al, 2022. AnoLLM: Large Language Models for Tabular Anomaly Detection.
\bibitem{rA415} HAN et al, 2024. HAN, J.; LIU, C.; GAO, P.; ZHOU, A.; HU, X.; YAN, S.; LU, P.; LI, H.; QIAO, Y. LLaMA-Adapter: Efficient Fine-tuning of Language Models with Zero-init Attention. ICLR, 2024. arXiv:2303.16199. DOI: https://doi.org/10.48550/arXiv.2303.16199 URL: https://arxiv.org/abs/2303.16199.
\bibitem{rA42} DETTMERS et al, 2023. Dettmers et al., 2023), além de variantes (AdaLoRA, DoRA, adapters, prefix-tuning, prompt-tuning, (IA)3, BitFit); (b) LLMs open-source candidatos a ajuste fino (LLaMA, Llama 2, Llama 3, Mistral); (c) fundamentos de few-shot / zero-shot / chain-of-thought que justificam o uso de LLMs sem ajuste de pesos; (d) adaptação de LLMs ao domínio de tráfego / telemetria de rede (NetGPT.
\bibitem{rA422} BROWN et al, 2005. BROWN, T. B.; MANN, B.; RYDER, N.; SUBBIAH, M.; KAPLAN, J.; DHARIWAL, P. et al. Language Models are Few-Shot Learners. NeurIPS, 2020. arXiv:2005.14165. DOI: https://doi.org/10.48550/arXiv.2005.14165 URL: https://arxiv.org/abs/2005.14165.
\bibitem{rA424} ACHIAM et al, 2023. ACHIAM, J.; ADLER, S.; AGARWAL, S.; AHMAD, L.; AKKAYA, I. et al. GPT-4 Technical Report. 2023. arXiv:2303.08774. DOI: https://doi.org/10.48550/arXiv.2303.08774 URL: https://arxiv.org/abs/2303.08774.
\bibitem{rA425} MENG et al, 2025. MENG, X.; LIN, C.; WANG, Y.; ZHANG, Y. NetGPT: Generative Pretrained Transformer for Network Traffic. 2023 (rev. 2025.
\bibitem{rA426} GINIGE et al, 2024. GINIGE, Y.; DAHANAYAKA, T.; HUANG, Y.; JOURJON, G.; SENEVIRATNE, S. TrafficGPT: An LLM Approach for Open-Set Encrypted Traffic Classification. Proceedings of the 19th Asian Internet Engineering Conference (AINTEC), 2024. DOI: https://doi.org/10.1145/3674213.3674217 URL: https://dl.acm.org/doi/10.1145/3674213.3674217.
\bibitem{rA427} CUI et al, 2025. CUI, T.; LIN, X.; LI, S.; CHEN, M.; YIN, Q.; LI, Q.; XU, K. TrafficLLM: Enhancing Large Language Models for Network Traffic Analysis with Generic Traffic Representation. 2025. arXiv:2504.04222. DOI: https://doi.org/10.48550/arXiv.2504.04222 URL: https://arxiv.org/abs/2504.04222.
\bibitem{rA43} ZHANG et al, 2023. When LLMs Meet Cybersecurity: A Systematic Literature Review.
\bibitem{rA431} SUI et al, 2024. SUI, Y.; ZHOU, M.; ZHOU, M.; HAN, S.; ZHANG, D. Table Meets LLM: Can Large Language Models Understand Structured Table Data? A Benchmark and Empirical Study. WSDM, 2024. arXiv:2305.13062. DOI: https://doi.org/10.48550/arXiv.2305.13062 URL: https://arxiv.org/abs/2305.13062.
\bibitem{rA432} HEGSELMANN et al, 2023. HEGSELMANN, S.; BUENDIA, A.; LANG, H.; AGRAWAL, M.; JIANG, X.; SONTAG, D. TabLLM: Few-shot Classification of Tabular Data with Large Language Models. AISTATS, 2023. arXiv:2210.10723. DOI: https://doi.org/10.48550/arXiv.2210.10723 URL: https://arxiv.org/abs/2210.10723.
\bibitem{rA44} LIU et al, 2024. On IoT intrusion detection based on data augmentation for enhancing learning on unbalanced samples.
\bibitem{rA45} HOULSBY et al, 1902. HOULSBY, N.; GIURGIU, A.; JASTRZEBSKI, S.; MORRONE, B.; DE LAROUSSILHE, Q.; GESMUNDO, A.; ATTARIYAN, M.; GELLY, S. Parameter-Efficient Transfer Learning for NLP. ICML, 2019. arXiv:1902.00751. DOI: https://doi.org/10.48550/arXiv.1902.00751 URL: https://arxiv.org/abs/1902.00751.
\bibitem{rA49} LIU et al, 2022. On IoT intrusion detection based on data augmentation for enhancing learning on unbalanced samples.
\bibitem{rA51} ESSENCIAL Danilevsky et al, 2010. Essencial / Recomendado / Complementar. - Estratégia de paralelização: busca distribuída em 3 eixos temáticos via subagentes (fundamentos de XAI; XAI para NIDS/HuntGPT; LLMs/calibração/ justificativa textual), consolidada e suplementada por buscas diretas do pesquisador principal. -------------------------------------------------------------------------------- 3. REFERÊNCIAS CLASSIFICADAS (.
\bibitem{rA511} ESSENCIAL Ali et al, 2023. Essencial / Recomendado / Complementar. - Estratégia de paralelização: busca distribuída em 3 eixos temáticos via subagentes (fundamentos de XAI; XAI para NIDS/HuntGPT; LLMs/calibração/ justificativa textual), consolidada e suplementada por buscas diretas do pesquisador principal. -------------------------------------------------------------------------------- 3. REFERÊNCIAS CLASSIFICADAS (.
\bibitem{rA512} ESSENCIAL Neupane et al, 2022. Essencial / Recomendado / Complementar. - Estratégia de paralelização: busca distribuída em 3 eixos temáticos via subagentes (fundamentos de XAI; XAI para NIDS/HuntGPT; LLMs/calibração/ justificativa textual), consolidada e suplementada por buscas diretas do pesquisador principal. -------------------------------------------------------------------------------- 3. REFERÊNCIAS CLASSIFICADAS (.
\bibitem{rA517} RECOMENDADO Sarker, 2024. Recomendado / Complementar. - Estratégia de paralelização: busca distribuída em 3 eixos temáticos via subagentes (fundamentos de XAI; XAI para NIDS/HuntGPT; LLMs/calibração/ justificativa textual), consolidada e suplementada por buscas diretas do pesquisador principal. -------------------------------------------------------------------------------- 3. REFERÊNCIAS CLASSIFICADAS (MÍNIMO 30 —.
\bibitem{rA52} ESSENCIAL Arrieta et al, 2019. Essencial / Recomendado / Complementar. - Estratégia de paralelização: busca distribuída em 3 eixos temáticos via subagentes (fundamentos de XAI; XAI para NIDS/HuntGPT; LLMs/calibração/ justificativa textual), consolidada e suplementada por buscas diretas do pesquisador principal. -------------------------------------------------------------------------------- 3. REFERÊNCIAS CLASSIFICADAS (.
\bibitem{rA520} COMPLEMENTAR Houssel et al, 2024. Complementar. - Estratégia de paralelização: busca distribuída em 3 eixos temáticos via subagentes (fundamentos de XAI; XAI para NIDS/HuntGPT; LLMs/calibração/ justificativa textual), consolidada e suplementada por buscas diretas do pesquisador principal. -------------------------------------------------------------------------------- 3. REFERÊNCIAS CLASSIFICADAS (MÍNIMO 30 — 38 VALIDADAS).
\bibitem{rA521} COMPLEMENTAR Pizarro et al, 2025. Complementar. - Estratégia de paralelização: busca distribuída em 3 eixos temáticos via subagentes (fundamentos de XAI; XAI para NIDS/HuntGPT; LLMs/calibração/ justificativa textual), consolidada e suplementada por buscas diretas do pesquisador principal. -------------------------------------------------------------------------------- 3. REFERÊNCIAS CLASSIFICADAS (MÍNIMO 30 — 38 VALIDADAS).
\bibitem{rA522} COMPLEMENTAR Lightweight Fine-Tuning of LLMs for Explainable Intrusion Detection in SDN. arXiv:2510.23313 (preprint), 2025. Complementar. - Estratégia de paralelização: busca distribuída em 3 eixos temáticos via subagentes (fundamentos de XAI; XAI para NIDS/HuntGPT; LLMs/calibração/ justificativa textual), consolidada e suplementada por buscas diretas do pesquisador principal. -------------------------------------------------------------------------------- 3. REFERÊNCIAS CLASSIFICADAS (MÍNIMO 30 — 38 VALIDADAS).
\bibitem{rA53} ESSENCIAL Ribeiro et al, 2016. Essencial / Recomendado / Complementar. - Estratégia de paralelização: busca distribuída em 3 eixos temáticos via subagentes (fundamentos de XAI; XAI para NIDS/HuntGPT; LLMs/calibração/ justificativa textual), consolidada e suplementada por buscas diretas do pesquisador principal. -------------------------------------------------------------------------------- 3. REFERÊNCIAS CLASSIFICADAS (.
\bibitem{rA531} ESSENCIAL Kadavath et al, 2022. Essencial / Recomendado / Complementar. - Estratégia de paralelização: busca distribuída em 3 eixos temáticos via subagentes (fundamentos de XAI; XAI para NIDS/HuntGPT; LLMs/calibração/ justificativa textual), consolidada e suplementada por buscas diretas do pesquisador principal. -------------------------------------------------------------------------------- 3. REFERÊNCIAS CLASSIFICADAS (.
\bibitem{rA532} RECOMENDADO Guo et al, 2017. Recomendado / Complementar. - Estratégia de paralelização: busca distribuída em 3 eixos temáticos via subagentes (fundamentos de XAI; XAI para NIDS/HuntGPT; LLMs/calibração/ justificativa textual), consolidada e suplementada por buscas diretas do pesquisador principal. -------------------------------------------------------------------------------- 3. REFERÊNCIAS CLASSIFICADAS (MÍNIMO 30 —.
\bibitem{rA533} RECOMENDADO Lin et al, 2022. Recomendado / Complementar. - Estratégia de paralelização: busca distribuída em 3 eixos temáticos via subagentes (fundamentos de XAI; XAI para NIDS/HuntGPT; LLMs/calibração/ justificativa textual), consolidada e suplementada por buscas diretas do pesquisador principal. -------------------------------------------------------------------------------- 3. REFERÊNCIAS CLASSIFICADAS (MÍNIMO 30 —.
\bibitem{rA534} RECOMENDADO Zhu et al, 2023. Recomendado / Complementar. - Estratégia de paralelização: busca distribuída em 3 eixos temáticos via subagentes (fundamentos de XAI; XAI para NIDS/HuntGPT; LLMs/calibração/ justificativa textual), consolidada e suplementada por buscas diretas do pesquisador principal. -------------------------------------------------------------------------------- 3. REFERÊNCIAS CLASSIFICADAS (MÍNIMO 30 —.
\bibitem{rA537} ESSENCIAL Wei et al, 2022. Essencial / Recomendado / Complementar. - Estratégia de paralelização: busca distribuída em 3 eixos temáticos via subagentes (fundamentos de XAI; XAI para NIDS/HuntGPT; LLMs/calibração/ justificativa textual), consolidada e suplementada por buscas diretas do pesquisador principal. -------------------------------------------------------------------------------- 3. REFERÊNCIAS CLASSIFICADAS (.
\bibitem{rA54} ESSENCIAL Lundberg et al, 2017. Essencial / Recomendado / Complementar. - Estratégia de paralelização: busca distribuída em 3 eixos temáticos via subagentes (fundamentos de XAI; XAI para NIDS/HuntGPT; LLMs/calibração/ justificativa textual), consolidada e suplementada por buscas diretas do pesquisador principal. -------------------------------------------------------------------------------- 3. REFERÊNCIAS CLASSIFICADAS (.
\bibitem{rA61} ESSENCIAL - Dataset UNSW-NB15 (descrição oficial / paper fundador) Moustafa et al, 2015. Essencial / Recomendado / Complementar. - Estratégia de paralelização: busca distribuída em 3 eixos temáticos via subagentes (fundamentos de XAI; XAI para NIDS/HuntGPT; LLMs/calibração/ justificativa textual), consolidada e suplementada por buscas diretas do pesquisador principal. -------------------------------------------------------------------------------- 3. REFERÊNCIAS CLASSIFICADAS (.
\bibitem{rA610} RECOMENDADO - Dataset InSDN et al, 2020. Recomendado / Complementar. - Estratégia de paralelização: busca distribuída em 3 eixos temáticos via subagentes (fundamentos de XAI; XAI para NIDS/HuntGPT; LLMs/calibração/ justificativa textual), consolidada e suplementada por buscas diretas do pesquisador principal. -------------------------------------------------------------------------------- 3. REFERÊNCIAS CLASSIFICADAS (MÍNIMO 30 —.
\bibitem{rA614} RECOMENDADO - Data augmentation em IoT (desbalanceamento) Zhang et al, 2022. Recomendado / Complementar. - Estratégia de paralelização: busca distribuída em 3 eixos temáticos via subagentes (fundamentos de XAI; XAI para NIDS/HuntGPT; LLMs/calibração/ justificativa textual), consolidada e suplementada por buscas diretas do pesquisador principal. -------------------------------------------------------------------------------- 3. REFERÊNCIAS CLASSIFICADAS (MÍNIMO 30 —.
\bibitem{rA615} RECOMENDADO - Geração sintética e privacy-preserving de traces Aceto et al, 2024. Recomendado / Complementar. - Estratégia de paralelização: busca distribuída em 3 eixos temáticos via subagentes (fundamentos de XAI; XAI para NIDS/HuntGPT; LLMs/calibração/ justificativa textual), consolidada e suplementada por buscas diretas do pesquisador principal. -------------------------------------------------------------------------------- 3. REFERÊNCIAS CLASSIFICADAS (MÍNIMO 30 —.
\bibitem{rA616} ESSENCIAL - AnoLLM et al, 2025. Essencial / Recomendado / Complementar. - Estratégia de paralelização: busca distribuída em 3 eixos temáticos via subagentes (fundamentos de XAI; XAI para NIDS/HuntGPT; LLMs/calibração/ justificativa textual), consolidada e suplementada por buscas diretas do pesquisador principal. -------------------------------------------------------------------------------- 3. REFERÊNCIAS CLASSIFICADAS (.
\bibitem{rA617} RECOMENDADO - Survey LLMs em dados tabulares Fang et al, 2024. Recomendado / Complementar. - Estratégia de paralelização: busca distribuída em 3 eixos temáticos via subagentes (fundamentos de XAI; XAI para NIDS/HuntGPT; LLMs/calibração/ justificativa textual), consolidada e suplementada por buscas diretas do pesquisador principal. -------------------------------------------------------------------------------- 3. REFERÊNCIAS CLASSIFICADAS (MÍNIMO 30 —.
\bibitem{rA618} RECOMENDADO - Detecção zero-shot de anomalias tabulares com LLMs Li et al, 2024. Recomendado / Complementar. - Estratégia de paralelização: busca distribuída em 3 eixos temáticos via subagentes (fundamentos de XAI; XAI para NIDS/HuntGPT; LLMs/calibração/ justificativa textual), consolidada e suplementada por buscas diretas do pesquisador principal. -------------------------------------------------------------------------------- 3. REFERÊNCIAS CLASSIFICADAS (MÍNIMO 30 —.
\bibitem{rA62} ESSENCIAL - Avaliação estatística do UNSW-NB15 Moustafa et al, 2016. Essencial / Recomendado / Complementar. - Estratégia de paralelização: busca distribuída em 3 eixos temáticos via subagentes (fundamentos de XAI; XAI para NIDS/HuntGPT; LLMs/calibração/ justificativa textual), consolidada e suplementada por buscas diretas do pesquisador principal. -------------------------------------------------------------------------------- 3. REFERÊNCIAS CLASSIFICADAS (.
\bibitem{rA623} RECOMENDADO - Augmentation GAN para classes minoritárias em rede Kang et al, 2025. Recomendado / Complementar. - Estratégia de paralelização: busca distribuída em 3 eixos temáticos via subagentes (fundamentos de XAI; XAI para NIDS/HuntGPT; LLMs/calibração/ justificativa textual), consolidada e suplementada por buscas diretas do pesquisador principal. -------------------------------------------------------------------------------- 3. REFERÊNCIAS CLASSIFICADAS (MÍNIMO 30 —.
\bibitem{rA629} COMPLEMENTAR - Testbed híbrido SDN et al, 2023. Complementar. - Estratégia de paralelização: busca distribuída em 3 eixos temáticos via subagentes (fundamentos de XAI; XAI para NIDS/HuntGPT; LLMs/calibração/ justificativa textual), consolidada e suplementada por buscas diretas do pesquisador principal. -------------------------------------------------------------------------------- 3. REFERÊNCIAS CLASSIFICADAS (MÍNIMO 30 — 38 VALIDADAS).
\bibitem{rA63} ESSENCIAL - Metodologia de geração CIC (base do CSE-CIC-IDS2018) Sharafaldin et al, 2018. Essencial / Recomendado / Complementar. - Estratégia de paralelização: busca distribuída em 3 eixos temáticos via subagentes (fundamentos de XAI; XAI para NIDS/HuntGPT; LLMs/calibração/ justificativa textual), consolidada e suplementada por buscas diretas do pesquisador principal. -------------------------------------------------------------------------------- 3. REFERÊNCIAS CLASSIFICADAS (.
\bibitem{rA630} RECOMENDADO - Survey de geração sintética de tráfego de rede Synthetic Network Traffic Data Generation: A Comparative Analysis. arXiv preprint, 2024. Recomendado / Complementar. - Estratégia de paralelização: busca distribuída em 3 eixos temáticos via subagentes (fundamentos de XAI; XAI para NIDS/HuntGPT; LLMs/calibração/ justificativa textual), consolidada e suplementada por buscas diretas do pesquisador principal. -------------------------------------------------------------------------------- 3. REFERÊNCIAS CLASSIFICADAS (MÍNIMO 30 —.
\bibitem{rA64} ESSENCIAL - Dataset CSE-CIC-IDS2018 (descrição oficial) Canadian Institute for Cybersecurity (CIC) \& Communications Security Establishment (CSE). CSE-CIC-IDS, 2018. Essencial / Recomendado / Complementar. - Estratégia de paralelização: busca distribuída em 3 eixos temáticos via subagentes (fundamentos de XAI; XAI para NIDS/HuntGPT; LLMs/calibração/ justificativa textual), consolidada e suplementada por buscas diretas do pesquisador principal. -------------------------------------------------------------------------------- 3. REFERÊNCIAS CLASSIFICADAS (.
\bibitem{rA65} ESSENCIAL - Survey de datasets de NIDS Ring et al, 2019. Essencial / Recomendado / Complementar. - Estratégia de paralelização: busca distribuída em 3 eixos temáticos via subagentes (fundamentos de XAI; XAI para NIDS/HuntGPT; LLMs/calibração/ justificativa textual), consolidada e suplementada por buscas diretas do pesquisador principal. -------------------------------------------------------------------------------- 3. REFERÊNCIAS CLASSIFICADAS (.
\bibitem{rA66} RECOMENDADO - Análise do KDD Cup 99 / NSL-KDD et al, 2009. Recomendado / Complementar. - Estratégia de paralelização: busca distribuída em 3 eixos temáticos via subagentes (fundamentos de XAI; XAI para NIDS/HuntGPT; LLMs/calibração/ justificativa textual), consolidada e suplementada por buscas diretas do pesquisador principal. -------------------------------------------------------------------------------- 3. REFERÊNCIAS CLASSIFICADAS (MÍNIMO 30 —.
\bibitem{rA69} ESSENCIAL - Mininet (emulador para SDN) Lantz et al, 2010. Essencial / Recomendado / Complementar. - Estratégia de paralelização: busca distribuída em 3 eixos temáticos via subagentes (fundamentos de XAI; XAI para NIDS/HuntGPT; LLMs/calibração/ justificativa textual), consolidada e suplementada por buscas diretas do pesquisador principal. -------------------------------------------------------------------------------- 3. REFERÊNCIAS CLASSIFICADAS (.
\bibitem{rRODRIGUES} Rodrigues, João Pedro Santos, 2026. Projeto de pesquisa: detecção e mitigação inteligente de anomalias em redes definidas por software com LLMs open-source.
\end{thebibliography}

\end{document}

